\begin{proposition}[Auxiliary torsor plumbing with cusp multiplicity twelve]\hypertarget{tag-030B}{}\label{push:src:B.3.7}\label{push:B37}
\StatementRecord{push:src:B.3.7}{proved}{030B}{\SeriesRef{push:src:B.3.6}, \SeriesRef{push:src:A.1.2}, \SeriesRef{push:src:A.2.2}; \cite[Proposition 0.9.2 and the opening
of Section 4.2, p.370]{CDL24}; \cite[Proposition 8.12, equation (8.25)]{Shioda90}}

For the torsor and affine lift $(V,A)$
of \ref{push:src:B.3.6}, there is a line bundle $M$ on
$V$ representing the narrow Jacobian
section $P$, with
\[
 M^2=-2,\qquad
 M\otimes(A^*M)^{-1}
       \simeq\pi^*\mathcal O_C(\infty).
\]
These data give a proper flat nodal
plumbing $\mathcal Z^\flat\to\Delta_h$
and two uniform smooth-total-space
semistable modifications
$\mathcal Z^\pm\to\Delta_h$, together
with holomorphic maps to $B$.
For $h\ne0$, the nodal fibre has
twelve ordinary double points and
Euler number twelve; each small-resolved
fibre has Euler number twenty-four.
The cusp fibre of the latter map to
$B$ has multiplicity twelve.
Consequently this auxiliary construction
does not fill the prescribed $S^6$
family. No proper birational modification
over $B$ can turn its multiple cusp
into a reduced cusp.
\end{proposition}


\begin{proof}
\begin{enumerate}[label=\arabic*.,leftmargin=*]
\item The section $P$ determines a relative
Picard class of multidegree zero on
every fibre component of $V\to C$.
It lifts to a line bundle $M$ because
$\operatorname{Br}(C)=0$; this is the
relative Picard consequence of
\cite[Proposition 0.9.2 and the opening
of Section 4.2, p.370]{CDL24}.
In local trivializations of the torsor,
translation acts trivially on the
relative degree-zero Picard class.
Thus the local Poincar\'e lines glue
up to base lines, with precisely that
Brauer obstruction. In particular we
may and do retain multidegree zero
on all twelve fibre components.

\item Choose a finite multisection cover of
degree $d$ which trivializes the torsor
generically, and resolve its base change.
There the pulled-back $M$ becomes the
narrow Poincar\'e line for the pulled-back
$P$, up to a base line. Its square is
$-2d$, by height base change
\cite[Proposition 8.12, equation (8.25)]{Shioda90}.
Base lines do not change the square
of a fibre-degree-zero line. The
intersection pullback formula therefore
gives $M^2=-2$.

\item Since $\sigma P=P$, the line
$J=M\otimes(A^*M)^{-1}$ is trivial
on every fibre, including the
$I_{12}$ fibre. Grauert base change
and the evaluation map give
$J=\pi^*J_C$ for a line on $C$.
The line $M\otimes(t_{6P}^*M)^{-1}$
has base degree twelve. To check the
sign and degree, pass to the same
trivializing cover. On a Jacobian,
the degree of the restriction of
$\mathcal O(P-O)$ to a section $Q$
minus its degree on $O$ is
$-\langle P,Q\rangle$, by the height
formula and the narrowness of $P$.
For $Q=6P$ this is $-12d$.
Descending the degree factor $d$
proves the assertion. Iterating $A^*$
twelve times now gives
$12\deg J_C=12$, so $\deg J_C=1$.

\item The forms $da/b^2$, $a\,da/b^2$ and
$da/b$ form a basis of $H^0(C,K_C)$.
Their $\sigma$-eigenvalues are
$\lambda^7,\lambda^{10},\lambda^{11}$.
Thus the characteristic polynomial
on $H^1(C,\mathbb Z)$ is
\[
 p(x)=\Phi_{12}(x)\Phi_6(x)
     =x^6-x^5+x^3-x+1,
 \qquad
 p(x)+(1-x)(x^5+x^2+x)=1.
\]
It follows that $1-\sigma^*$ is an
automorphism of $\operatorname{Pic}^0(C)$,
with inverse $(\sigma^*)^5+(\sigma^*)^2+\sigma^*$.
Twisting $M$ by a base line changes
$J_C$ by $N\otimes(\sigma^*N)^{-1}$.
The displayed inverse therefore supplies
a base twist making
$J_C=\mathcal O_C(\infty)$.
Let $\phi$ be its section with divisor
the reduced $I_{12}$ fibre $F_\infty$.

\item Apply the plumbing maps of \ref{push:src:A.1.2}
to the data $(V,M,A)$.
Set $R_V=\mathbb P_V(\mathcal O\oplus M)$,
and identify its zero and infinity
sections by $s_0(x)\sim s_\infty(Ax)$.
On the node neighbourhood over $V$ use
\[
 a\in M_x,\quad b\in M_{Ax}^{-1},
 \qquad a\otimes b=h\phi_x.
\]
The $a$ collar is glued to the annular
core at $(x,a)$; the $b$ collar is glued
at $(Ax,b^{-1})$. Choose the Hermitian
metric and collar radii as in \ref{push:src:A.1.2}.
Compactness of $V$ and the bounded norm
of $\phi$ give disjoint collars,
Hausdorffness and compact representatives
by exactly the same inequalities.
These are holomorphic gluing maps of
fourfolds over $\Delta_h$.
Their local hypersurface rings are flat
over $h$, and the compact representatives
prove properness. The map to $B$ is
well defined because $t\circ\sigma=t$.

\item At a smooth point of $F_\infty$ the
equation is $ab=hz_C$; at one of its
twelve nodes, $z_C=uv$ and it is
$ab=huv$. For $h\ne0$ the latter
has one ordinary double point.
These twelve centres are distinct
even if $A$ rotates the fibre
components. A centre has $a=b=0$
and is in neither annular collar;
no overlap identifies it with the
centre indexed by $Ax$.
Thus the double locus is parametrized by $V$.
The two uniform prism subdivisions
of \ref{push:src:A.2.2} extend over these charts:
the affine action preserves the
orientation of the $I_{12}$ cycle,
and all remaining changes are units.
They give proper small modifications,
smooth total spaces and reduced
normal-crossings $h$-central fibres.
Equivalently their normalizations are
the blow-up of $R_V$ along $F_\infty$
in one section, followed by the twelve
chosen small resolutions. The normal
product becomes
$\pi^*\mathcal O_C(\infty)(-F_\infty)=\mathcal O_V$.

\item The nearby function of the nodal
plumbing is one on $F_\infty$ and
zero on its complement in the conductor.
Since $e(F_\infty)=12$, its integral
is twelve. Each small resolution
replaces a point by $\mathbb P^1$
and adds one, giving twenty-four.
Independently $e(R_V)=24$; above each
of the twelve nodes the modified
normalization replaces a point by
a chain of two projective lines,
of Euler number three. The open
parts of the fibre are copies of
$\mathbb C^*$ and contribute zero.
The modified normalization therefore
has Euler number $48$, its pinching
has Euler number $36$, and the
semistable nearby count is
$48-2e(V)=24$.

\item Finally $z=a/b$ is a coordinate on
$C$ near infinity and
$q_B=1/t=z^{12}/c^3$, with $c$ a unit.
At a generic point of the auxiliary
cusp fibre this gives multiplicity
twelve. In fact the map $q_B$ is
locally a twelfth power throughout
the smooth nearby threefold; a
holomorphic root of the unit can
be chosen locally. On any common
resolution of a proper birational
comparison over $B$, every divisor
over that fibre consequently has
$q_B$-valuation divisible by twelve.
A reduced fibre on another smooth
model would have a strict transform
of valuation one, a contradiction.

\end{enumerate}
\end{proof}
